\documentclass[11pt,a4paper]{article}

% ----------------------------------------------------------------------------
% Protokoll zur Projektarbeit "Maschinelles Sehen"
% Pflanzen-, Krankheits- und Zimmerpflanzenerkennung via Transfer Learning
% ----------------------------------------------------------------------------

\usepackage[utf8]{inputenc}
\usepackage[T1]{fontenc}
\usepackage[ngerman]{babel}
\usepackage{lmodern}
\usepackage{microtype}
\usepackage{graphicx}
\usepackage{booktabs}
\usepackage{array}
\usepackage{tabularx}
\usepackage{xcolor}
\usepackage{listings}
\usepackage{amsmath}
\usepackage{amssymb}
\usepackage{textcomp}
\usepackage{enumitem}
\usepackage[hidelinks]{hyperref}
\usepackage{caption}
\usepackage{geometry}
\geometry{margin=2.5cm}

% Farben fuer Code-Listings
\definecolor{codebg}{rgb}{0.96,0.96,0.96}
\definecolor{codekw}{rgb}{0.13,0.29,0.53}
\definecolor{codecm}{rgb}{0.40,0.55,0.40}
\definecolor{codestr}{rgb}{0.58,0.20,0.20}

\lstset{
  basicstyle=\ttfamily\footnotesize,
  backgroundcolor=\color{codebg},
  keywordstyle=\color{codekw}\bfseries,
  commentstyle=\color{codecm}\itshape,
  stringstyle=\color{codestr},
  numbers=left, numberstyle=\tiny\color{gray}, numbersep=8pt,
  frame=single, rulecolor=\color{gray!40},
  breaklines=true, showstringspaces=false,
  language=Python, captionpos=b, xleftmargin=1.5em,
}

% Hilfsmakro: Bild nur einbinden, wenn vorhanden (Protokoll kompiliert sonst trotzdem)
\newcommand{\maybegraphic}[2]{%
  \IfFileExists{#1}{\includegraphics[#2]{#1}}{%
    \fbox{\parbox[c][3cm][c]{0.6\linewidth}{\centering\itshape Abbildung nicht gefunden:\\\texttt{#1}}}}}

\title{\textbf{Pflanzen-, Krankheits- und Zimmerpflanzenerkennung\\mit Transfer Learning}\\[0.4em]
\large Protokoll zur Projektarbeit im Modul \emph{Maschinelles Sehen}}
\author{Kevin Ruppaner}
\date{\today}

\begin{document}
\maketitle

\begin{abstract}
\noindent
In dieser Projektarbeit wird ein vortrainiertes neuronales Netz (CNN) mittels
\emph{Transfer Learning} darauf angepasst, Pflanzen auf Fotos zu erkennen. Es
entstanden drei Aufgabenstellungen: (1)~die Erkennung von \textbf{Nutzpflanzen
und deren Krankheiten} aus Blattfotos (38~Klassen, PlantVillage), (2)~die
Bestimmung der \textbf{Zimmerpflanzen-Art} (47~Klassen) und (3)~eine binäre
\textbf{Zustandsbewertung} (gesund vs.\ welk). Auf den Validierungsdaten werden
sehr hohe Genauigkeiten erreicht (Krankheiten bis 99{,}7\,\%, Arten 91{,}3\,\%).
Ein abschließender Test mit zehn \emph{unabhängigen} Fotos aus dem Internet
ordnet die Ergebnisse ehrlich ein: hier wurden 8 von 10 Arten korrekt erkannt.
Das Protokoll erläutert Idee, Aufbau, Methodik und Ergebnisse und diskutiert
Grenzen sowie konkrete Verbesserungsmöglichkeiten.
\end{abstract}

\tableofcontents
\bigskip

% ============================================================================
\section{Idee und Zielsetzung}
% ============================================================================

\subsection{Motivation}
Pflanzenkrankheiten früh zu erkennen ist in der Landwirtschaft wirtschaftlich
bedeutsam, und auch im Hobbybereich (Zimmerpflanzen) ist die Frage „Welche
Pflanze ist das und geht es ihr gut?“ alltagsnah. Beide Probleme lassen sich als
\textbf{Bildklassifikation} formulieren: Eingabe ist ein Foto, Ausgabe ist eine
Klassenbezeichnung. Genau dafür sind \emph{Convolutional Neural Networks} (CNNs)
das Standardwerkzeug des maschinellen Sehens.

\subsection{Kernidee: Transfer Learning}
Ein CNN von Grund auf zu trainieren benötigt sehr viele Bilder und viel
Rechenzeit. Stattdessen nutzen wir \textbf{Transfer Learning}: Wir nehmen ein
Netz, das bereits auf dem großen Bilddatensatz \emph{ImageNet}
(1{,}2~Mio.\ Bilder, 1000~Klassen) trainiert wurde. Dieses Netz hat allgemeine
visuelle Merkmale (Kanten, Texturen, Formen) bereits gelernt. Wir tauschen nur
den letzten Klassifikationskopf aus und passen das Netz an unsere Pflanzendaten
an. So genügen vergleichsweise wenige Bilder und Minuten statt Tage Trainingszeit
-- ideal für eine Projektarbeit.

\subsection{Lernziel der Arbeit}
Da es sich um eine Lern\-/Übungsarbeit handelt, liegt der Fokus nicht auf einem
produktionsreifen System, sondern auf dem \emph{Verständnis} der Pipeline:
Datenaufbereitung $\rightarrow$ Modellaufbau $\rightarrow$ Zwei-Phasen-Training
$\rightarrow$ Evaluation $\rightarrow$ Inferenz auf neuen Bildern, sowie dem
ehrlichen Einordnen der Ergebnisse.

% ============================================================================
\section{Theoretischer Hintergrund}
% ============================================================================

\subsection{Convolutional Neural Networks (CNN)}
Ein CNN verarbeitet ein Bild über mehrere \emph{Faltungsschichten}
(Convolutions), die lokale Muster erkennen, gefolgt von Pooling-Schichten, die
die Auflösung reduzieren. Frühe Schichten reagieren auf einfache Merkmale
(Kanten), tiefere auf komplexe (z.\,B.\ Blattadern, Flecken). Am Ende steht ein
\emph{Klassifikationskopf} (eine voll verbundene Schicht), der die gelernten
Merkmale auf die Zielklassen abbildet.

\subsection{Verwendete Backbones}
Als „Backbone“ (Merkmalsextraktor) wurden zwei etablierte Architekturen aus
\texttt{torchvision} verglichen:
\begin{itemize}[nosep]
  \item \textbf{ResNet18} -- 11{,}2~Mio.\ Parameter, schnell, einsteigerfreundlich
        (Residual-Verbindungen erleichtern das Training tiefer Netze).
  \item \textbf{EfficientNet-B0} -- 4{,}1~Mio.\ Parameter, modernere Architektur,
        oft etwas genauer, aber langsamer im Training.
\end{itemize}

\subsection{Zwei-Phasen-Training}
\label{sec:zweiphasen}
Das Anpassen erfolgt in zwei Phasen -- ein in der Praxis bewährtes Vorgehen:
\begin{enumerate}[nosep]
  \item \textbf{Feature Extraction:} Das Backbone wird „eingefroren“
        (Gewichte unveränderlich), nur der neue Kopf lernt. Das ist schnell und
        stabil, weil nur wenige Parameter angepasst werden.
  \item \textbf{Fine-Tuning:} Das gesamte Netz wird mit einer \emph{kleinen}
        Lernrate weitertrainiert. So passen sich auch die tieferen Schichten
        sanft an die Pflanzendomäne an, ohne das Vorwissen zu „vergessen“.
\end{enumerate}
Im Fine-Tuning kommt zusätzlich ein \emph{Cosine-Annealing}-Scheduler zum
Einsatz, der die Lernrate über die Epochen langsam absenkt.

\subsection{Normalisierung und Augmentierung}
Die Bilder werden auf $224\times224$ skaliert und mit der ImageNet-Statistik
normalisiert (Mittelwert/Std je Farbkanal) -- genau so, wie die vortrainierten
Backbones es erwarten. Beim Training wird zusätzlich \emph{augmentiert}
(zufälliges Zuschneiden, Spiegeln, Rotieren, Farbänderung), damit das Modell
robuster wird und nicht einfach Trainingsbilder auswendig lernt.

\subsection{Bewertungsmetriken: Precision, Recall, F1}
\label{sec:metriken}
Die \emph{Accuracy} (Anteil korrekter Vorhersagen) allein ist trügerisch, wenn
die Klassen unausgewogen sind. Aussagekräftiger sind drei Maße, die sich aus den
vier Grundfällen einer Klasse ableiten: richtig-positiv (TP), falsch-positiv
(FP), falsch-negativ (FN), richtig-negativ (TN).
\begin{align*}
\text{Precision} &= \frac{TP}{TP + FP}
   && \text{\small(Wie viele als X erkannte sind wirklich X?)}\\[2pt]
\text{Recall} &= \frac{TP}{TP + FN}
   && \text{\small(Wie viele echte X wurden gefunden?)}\\[2pt]
\text{F1} &= 2\cdot\frac{\text{Precision}\cdot\text{Recall}}{\text{Precision}+\text{Recall}}
   && \text{\small(harmonisches Mittel beider)}
\end{align*}
Der \textbf{F1-Score} ist hoch nur dann, wenn \emph{beide} Werte hoch sind, und
bestraft so einseitige Modelle. Über mehrere Klassen aggregiert man entweder als
\textbf{Macro-F1} (ungewichteter Mittelwert -- jede Klasse zählt gleich, auch
seltene) oder als \textbf{Weighted-F1} (nach Klassengröße gewichtet). In dieser
Arbeit wird der Macro-F1 berichtet, weil er seltene Klassen nicht „verschwinden“
lässt.

\subsection{Konfusionsmatrix}
\label{sec:konfusion}
Eine \textbf{Konfusionsmatrix} (engl.\ \emph{confusion matrix}) stellt
Verwechslungen tabellarisch dar: Zeilen sind die \emph{wahren} Klassen, Spalten
die \emph{vorhergesagten}. Auf der Diagonale stehen die korrekten Fälle; alles
außerhalb sind Verwechslungen. Eine „saubere“ Diagonale bedeutet ein gutes
Modell, helle Nebenfelder zeigen, \emph{welche} Klassen verwechselt werden. Aus
ihr lassen sich Precision und Recall direkt ablesen. Die Auswertung wird im
Projekt von \texttt{src/evaluate.py} erzeugt:
\begin{lstlisting}[caption={Konfusionsmatrix + Report (\texttt{src/evaluate.py})}]
from sklearn.metrics import classification_report, confusion_matrix
print(classification_report(y_true, y_pred, target_names=class_names))
cm = confusion_matrix(y_true, y_pred)   # Zeilen=wahr, Spalten=vorhergesagt
\end{lstlisting}

% ============================================================================
\section{Projektaufbau}
% ============================================================================

\subsection{Verzeichnisstruktur}
Das Projekt trennt sauber zwischen wiederverwendbarem Quellcode (\texttt{src/}),
einem Benchmark-Werkzeug (\texttt{bench/}), einer Demo-App (\texttt{app/}) und den
Ergebnissen (\texttt{results/}, \texttt{models/}).

\begin{lstlisting}[language=bash, numbers=none, caption={Projektstruktur (ohne Daten/Modelle)}]
src/data.py             Transforms + DataLoader (ImageFolder), Label-Split
src/model.py            build_model(): Backbone + neuer Kopf, Checkpoint speichern/laden
src/train.py            Zwei-Phasen-Trainingsschleife, Geraeteerkennung (GPU/CPU)
src/evaluate.py         Accuracy, Classification-Report, Confusion Matrix
src/predict.py          Einzelbild-Vorhersage  -> (Pflanze, Krankheit, Konfidenz)
src/predict_houseplant.py  Kombinierte Inferenz: Art-Modell -> Zustand-Modell
bench/benchmark.py      Einheitlicher Trainings-/Evaluations-/Testlauf je Modell
bench/compare.py        Fasst alle Laeufe zu results/BENCHMARK.md zusammen
bench/plot.py           Diagramme (Accuracy vs. Zeit, Lernkurven)
app/app.py              Gradio-Weboberflaeche (lokal, CPU)
notebooks/              Trainings-Notebook
\end{lstlisting}

\subsection{Datenpipeline (\texttt{src/data.py})}
Die Daten liegen als \emph{ImageFolder} vor: ein Unterordner pro Klasse, der
Ordnername ist das Label. Zwei Layouts werden automatisch unterstützt:
ein bereits geteilter Datensatz (\texttt{train/} + \texttt{valid/}) oder ein
einzelner Ordner, der dann mit festem Zufalls-Seed in 80\,\% Training und
20\,\% Validierung geteilt wird. Für das stark unausgewogene Arten-Dataset
existiert optional ein \texttt{WeightedRandomSampler}, der seltene Klassen
häufiger zieht.

Ein kleiner, aber didaktisch schöner Kniff ist die Zerlegung des
PlantVillage-Labels in Pflanze und Zustand:
\begin{lstlisting}[caption={Label-Split in Pflanze und Krankheit (\texttt{src/data.py})}]
def split_label(label: str) -> tuple[str, str]:
    # "Tomato___Late_blight" -> ("Tomato", "Late blight")
    if "___" in label:
        plant, disease = label.split("___", 1)
    else:
        plant, disease = label, "unknown"
    return plant.replace("_", " ").strip(), disease.replace("_", " ").strip()
\end{lstlisting}

\subsection{Modellaufbau (\texttt{src/model.py})}
Die Funktion \texttt{build\_model()} lädt das gewählte Backbone mit
ImageNet-Gewichten, friert es bei Bedarf ein und ersetzt den Kopf durch eine
neue lineare Schicht mit \texttt{num\_classes} Ausgängen. Beim Speichern werden
neben den Gewichten auch die Klassennamen und der Backbone-Typ im Checkpoint
abgelegt, sodass die Inferenz das Modell vollständig rekonstruieren kann.

\subsection{Geräteerkennung}
Das Training erkennt die Hardware automatisch (CUDA / AMD-ROCm / DirectML /
Apple-MPS / CPU). Lokal wurde auf einer \textbf{AMD Radeon RX 5700 XT} unter
Windows via \texttt{torch-directml} trainiert; die Inferenz läuft auf der CPU.

% ============================================================================
\section{Datensätze}
% ============================================================================

\begin{table}[h]
\centering
\caption{Verwendete Datensätze}
\label{tab:datasets}
\small
\begin{tabularx}{\linewidth}{@{}lXrr@{}}
\toprule
\textbf{Aufgabe} & \textbf{Datensatz} & \textbf{Klassen} & \textbf{Bilder} \\
\midrule
Krankheiten (crop\_disease)      & PlantVillage (New Plant Diseases, augm.) & 38 & 87.867 \\
Zimmerpflanzen-Art               & House Plant Species                      & 47 & 14.774 \\
Zustand (gesund/welk)            & Healthy vs.\ Wilted Houseplants          &  2 &    904 \\
\bottomrule
\end{tabularx}
\end{table}

\noindent
\textbf{Wichtige Einschränkung:} PlantVillage besteht aus \emph{Laboraufnahmen}
(einzelnes Blatt, gleichmäßiger Hintergrund). Darauf werden sehr hohe
Genauigkeiten erreicht, die aber nur eingeschränkt auf reale Feldfotos
übertragbar sind (siehe Abschnitt~\ref{sec:diskussion}).

\subsection{Zwei-Modell-Architektur für Zimmerpflanzen}
Für Zimmerpflanzen existieren kaum Krankheitsdaten. Statt eines kombinierten
\texttt{Art\_\_\_Zustand}-Klassifikators wurden daher bewusst \textbf{zwei
getrennte Modelle} trainiert: ein datenreiches \emph{Art-Modell} (47~Arten) und
ein datenarmes binäres \emph{Zustand-Modell}. Bei der Inferenz wird ein Foto
nacheinander durch beide Modelle geschickt (siehe \texttt{predict\_houseplant.py}).

% ============================================================================
\section{Durchführung}
% ============================================================================

\subsection{Benchmark-Harness}
Damit alle Modelle fair und reproduzierbar verglichen werden, kapselt
\texttt{bench/benchmark.py} einen kompletten Lauf: Training (zwei Phasen),
Evaluation (Accuracy, Per-Klassen-Report, Confusion Matrix), Messung der
Inferenzgeschwindigkeit und Test auf unabhängigen Bildern. Pro Lauf entstehen
maschinenlesbare \texttt{metrics.json} sowie menschenlesbare Protokolle. Eine
\emph{append-only}-Historie (\texttt{results/history.jsonl}) hält jeden Lauf
dauerhaft fest, sodass der Effekt von Hyperparameter-Änderungen über die Zeit
nachvollziehbar bleibt.

\begin{lstlisting}[language=bash, numbers=none, caption={Beispielaufrufe}]
# Krankheiten, ResNet18
python -m bench.benchmark --backbone resnet18

# Zimmerpflanzen-Arten
python -m bench.benchmark --backbone resnet18 --data-root data/houseplants/house_plant_species \
    --run-name houseplant_species --task houseplant_species

# Vergleichstabelle + Diagramme erzeugen
python -m bench.compare && python -m bench.plot
\end{lstlisting}

\subsection{Hyperparameter}
Alle Läufe nutzten dieselben Einstellungen, um vergleichbar zu bleiben:
2~Kopf-Epochen + 3~Fine-Tuning-Epochen, Batch-Größe~32, Lernrate
$10^{-3}$ (Kopf) bzw.\ $10^{-4}$ (Fine-Tuning), Bildgröße~$224$, Optimierer
Adam, Verlustfunktion Cross-Entropy.

% ============================================================================
\section{Ergebnisse}
% ============================================================================

\subsection{Krankheitserkennung: ResNet18 vs.\ EfficientNet-B0}
\begin{table}[h]
\centering
\caption{Modellvergleich auf PlantVillage (38 Klassen, Validierungssplit)}
\label{tab:crop}
\small
\begin{tabular}{@{}lrrrrrr@{}}
\toprule
\textbf{Modell} & \textbf{Val-Acc} & \textbf{Macro-F1} & \textbf{Params} & \textbf{Train-Zeit} & \textbf{Inferenz} & \textbf{Testbilder} \\
\midrule
EfficientNet-B0 & 99{,}69\,\% & 0{,}9969 & 4{,}1\,M & 90\,min & 2{,}89\,ms & 32/33 \\
ResNet18        & 99{,}52\,\% & 0{,}9951 & 11{,}2\,M & 25\,min & 0{,}54\,ms & 33/33 \\
\bottomrule
\end{tabular}
\end{table}

\noindent
\textbf{Beobachtung:} EfficientNet-B0 ist minimal genauer, benötigt aber
$\sim$3{,}6$\times$ mehr Trainingszeit und ist in der Inferenz $\sim$5$\times$
langsamer. Für dieses Projekt ist \textbf{ResNet18} der bessere Kompromiss --
nahezu gleiche Genauigkeit bei deutlich geringerem Aufwand. Beide verwechseln vor
allem schwer unterscheidbare Mais-Krankheiten (Cercospora vs.\ Northern Leaf
Blight).

\subsection{Zimmerpflanzen: Arten und Zustand}
\begin{table}[h]
\centering
\caption{Zimmerpflanzen-Modelle (ResNet18)}
\label{tab:house}
\small
\begin{tabular}{@{}lrrrr@{}}
\toprule
\textbf{Modell} & \textbf{Klassen} & \textbf{Val-Acc} & \textbf{Macro-F1} & \textbf{Train-Zeit} \\
\midrule
Art (House Plant Species)        & 47 & 91{,}33\,\% & 0{,}9038 & 13\,min \\
Art + WeightedRandomSampler      & 47 & 90{,}93\,\% & 0{,}9038 & 12\,min \\
Zustand (gesund/welk)            &  2 & 88{,}33\,\% & 0{,}8833 &  2{,}5\,min \\
\bottomrule
\end{tabular}
\end{table}

\noindent
Das Art-Modell erreicht solide 91\,\% bei 47~Klassen. Schwächste Klassen sind
\emph{Yucca} (F1~0{,}53) und \emph{Begonia} -- teils wegen weniger Trainingsbildern,
teils wegen optischer Ähnlichkeit. Der \texttt{WeightedRandomSampler} brachte hier
keinen messbaren Vorteil (gleiches Macro-F1). Das Zustand-Modell ist mit nur
904~Bildern erwartungsgemäß schwächer.

\subsection{Lernkurven}
\begin{figure}[h]
\centering
\maybegraphic{../results/crop_disease_val_acc_curves.png}{width=0.48\linewidth}
\hfill
\maybegraphic{../results/houseplant_species_val_acc_curves.png}{width=0.48\linewidth}
\caption{Validierungs-Accuracy je Epoche: Krankheiten (links), Arten (rechts).
Der deutliche Sprung markiert jeweils den Übergang von Feature-Extraction zu
Fine-Tuning.}
\end{figure}

\subsection{F1-Auswertung und Konfusionsmatrizen}
\label{sec:f1konfusion}
Tabelle~\ref{tab:f1} fasst die F1-Werte aller Modelle zusammen (vgl.\ die
Definition in Abschnitt~\ref{sec:metriken}). Dass Macro-F1 und Accuracy bei den
Krankheiten fast identisch sind, liegt an den ausgewogenen Klassen; beim
Arten-Modell ist der Macro-F1 (0{,}904) spürbar niedriger als die Accuracy
(0{,}913), weil einige seltene Arten den ungewichteten Mittelwert drücken.

\begin{table}[h]
\centering
\caption{Precision, Recall und F1 (Macro-Mittel über alle Klassen, Validierung)}
\label{tab:f1}
\small
\begin{tabular}{@{}lrrrr@{}}
\toprule
\textbf{Modell} & \textbf{Accuracy} & \textbf{Macro-Precision} & \textbf{Macro-Recall} & \textbf{Macro-F1} \\
\midrule
EfficientNet-B0 (Krankheiten) & 99{,}69\,\% & 0{,}9970 & 0{,}9968 & 0{,}9969 \\
ResNet18 (Krankheiten)        & 99{,}52\,\% & 0{,}9952 & 0{,}9950 & 0{,}9951 \\
Art (47 Klassen)              & 91{,}33\,\% & 0{,}9100 & 0{,}9020 & 0{,}9038 \\
Zustand (gesund/welk)         & 88{,}33\,\% & 0{,}8834 & 0{,}8833 & 0{,}8833 \\
\bottomrule
\end{tabular}
\end{table}

\paragraph{Schwächste Klassen im Backbone-Vergleich.}
Der Macro-F1 verbirgt, \emph{wo} ein Modell schwächelt. Tabelle~\ref{tab:worstf1}
listet die drei Klassen mit dem niedrigsten F1 für beide Krankheits-Backbones.
Beide kämpfen mit denselben Klassen -- vor allem zwei optisch sehr ähnlichen
Mais-Krankheiten. EfficientNet-B0 liegt dort durchgängig etwas höher als
ResNet18, was den minimalen Gesamtvorsprung erklärt; der Unterschied ist aber so
klein, dass er die deutlich längere Trainingszeit kaum rechtfertigt.

\begin{table}[h]
\centering
\caption{F1-Score der drei schwächsten Klassen je Backbone (Krankheiten)}
\label{tab:worstf1}
\small
\begin{tabular}{@{}lcc@{}}
\toprule
\textbf{Klasse} & \textbf{ResNet18} & \textbf{EfficientNet-B0} \\
\midrule
Corn -- Cercospora / Gray leaf spot & 0{,}957 & 0{,}972 \\
Corn -- Northern Leaf Blight        & 0{,}962 & 0{,}975 \\
Tomato -- Target Spot               & 0{,}979 & 0{,}987 \\
\bottomrule
\end{tabular}
\end{table}

\paragraph{Konfusionsmatrix am einfachen Beispiel.}
Das binäre Zustand-Modell eignet sich, um eine Konfusionsmatrix vollständig
„von Hand“ zu lesen (je 90 Validierungsbilder pro Klasse):

\begin{table}[h]
\centering
\caption{Konfusionsmatrix Zustand-Modell (Zeilen = wahr, Spalten = vorhergesagt)}
\label{tab:cm2}
\small
\begin{tabular}{@{}l|cc|c@{}}
\toprule
 & \textbf{pred.\ gesund} & \textbf{pred.\ welk} & \textbf{Recall} \\
\midrule
\textbf{wahr gesund} & 80 (TP) & 10 (FN) & 0{,}889 \\
\textbf{wahr welk}   & 11 (FP) & 79 (TN) & 0{,}878 \\
\midrule
\textbf{Precision}   & 0{,}879 & 0{,}888 & \\
\bottomrule
\end{tabular}
\end{table}

\noindent
Aus der Matrix folgt für die Klasse „gesund“ direkt:
$\text{Precision}=\tfrac{80}{80+11}=0{,}879$,
$\text{Recall}=\tfrac{80}{80+10}=0{,}889$ und damit
$\text{F1}=2\cdot\tfrac{0{,}879\cdot0{,}889}{0{,}879+0{,}889}=0{,}884$.
Die Verwechslungen sind nahezu symmetrisch -- das Modell hat keine systematische
Schlagseite, ist aber mit nur 904~Trainingsbildern insgesamt noch unsicher.

\paragraph{Konfusionsmatrizen der Mehrklassen-Modelle.}
Bei 38 bzw.\ 47 Klassen ist die Matrix zu groß für eine Tabelle und wird als
Heatmap dargestellt (normalisiert je Zeile, dunkler = höher). Die stark
ausgeprägte Diagonale bestätigt die hohe Genauigkeit; die wenigen helleren
Nebenfelder entsprechen den in Abschnitt~\ref{sec:diskussion} genannten
Verwechslungen (z.\,B.\ ähnliche Mais-Krankheiten bzw.\ Yucca/Begonia).

\begin{figure}[h]
\centering
\maybegraphic{../results/resnet18/confusion_matrix.png}{width=0.49\linewidth}
\hfill
\maybegraphic{../results/houseplant_species/confusion_matrix.png}{width=0.49\linewidth}
\caption{Normalisierte Konfusionsmatrizen: ResNet18 auf den Krankheiten
(38~Klassen, links) und das Arten-Modell (47~Klassen, rechts). Die klare
Diagonale zeigt, dass die meisten Klassen sicher getrennt werden.}
\end{figure}

% ============================================================================
\section{Unabhängiger Test mit 10 Internet-Fotos}
% ============================================================================
\label{sec:webtest}

Die Validierungs-Accuracy stammt aus demselben Datensatz wie das Training und ist
daher optimistisch. Für eine ehrlichere Einschätzung wurden \textbf{zehn fremde
Fotos} (Wikimedia/Wikipedia, reale Aufnahmen, keine Laborbilder) durch die
kombinierte Inferenz \texttt{Art~$\rightarrow$~Zustand} geschickt.

\begin{table}[h]
\centering
\caption{Vorhersagen auf 10 unabhängigen Internet-Fotos}
\label{tab:webtest}
\small
\begin{tabular}{@{}llrlr@{}}
\toprule
\textbf{Foto} & \textbf{Vorhergesagte Art} & \textbf{Konf.} & \textbf{Zustand} & \textbf{Treffer} \\
\midrule
Aloe Vera        & Aloe Vera                    & 90{,}8\,\% & welk\,(56\%)   & \checkmark \\
Monstera         & Monstera Deliciosa           & 99{,}9\,\% & gesund\,(94\%) & \checkmark \\
Bogenhanf        & Snake plant (Sanseviera)     & 87{,}2\,\% & gesund\,(52\%) & \checkmark \\
Orchidee         & Orchid                       & 53{,}1\,\% & gesund\,(73\%) & \checkmark \\
Friedenslilie    & Peace lily                   & 99{,}8\,\% & gesund\,(97\%) & \checkmark \\
Gummibaum        & \emph{Schefflera}            & 48{,}2\,\% & gesund\,(95\%) & \texttimes \\
Efeutute (Pothos)& \emph{Elephant Ear (Alocasia)} & 36{,}9\,\% & gesund\,(51\%) & \texttimes \\
ZZ-Pflanze       & ZZ Plant (Zamioculcas)       & 100\,\%    & gesund\,(99\%) & \checkmark \\
Geldbaum (Jade)  & Jade plant (Crassula ovata)  & 94{,}0\,\% & gesund\,(93\%) & \checkmark \\
Tulpe            & Tulip                        & 99{,}9\,\% & gesund\,(83\%) & \checkmark \\
\midrule
\multicolumn{4}{r}{\textbf{Art korrekt:}} & \textbf{8 / 10} \\
\bottomrule
\end{tabular}
\end{table}

\noindent
\textbf{Interpretation:} 8 von 10 Arten wurden korrekt erkannt -- ein
realistisches Ergebnis, das deutlich unter den 91\,\% Validierungs-Accuracy
liegt und damit den Effekt des \emph{Domänenunterschieds} (andere Hintergründe,
Beleuchtung, Bildausschnitte) zeigt. Die beiden Fehler sind nachvollziehbar:
Der Gummibaum (glänzende, gegenständige Blätter) wurde mit \emph{Schefflera}
verwechselt, die Efeutute mit \emph{Elephant Ear} -- beide haben ähnlich
herzförmige Blätter. Auffällig sind die \emph{niedrigen Konfidenzen} (48\,\%,
37\,\%) bei den Fehlern: Das Modell ist sich hier selbst unsicher, was praktisch
nutzbar wäre (z.\,B.\ „unsicher“ ausgeben statt falscher Sicherheit).

Das Zustand-Modell ist erkennbar schwächer kalibriert: gesunde Pflanzen wurden
teils nur knapp als „gesund“ eingestuft (Bogenhanf 52\,\%) oder fälschlich als
„welk“ (Aloe 56\,\%). Mit nur 904~Trainingsbildern ist das nicht überraschend.

\begin{figure}[h]
\centering
\maybegraphic{../data/web_test/monstera.jpg}{width=0.3\linewidth}
\hfill
\maybegraphic{../data/web_test/rubber_plant.jpg}{width=0.3\linewidth}
\hfill
\maybegraphic{../data/web_test/pothos.jpg}{width=0.3\linewidth}
\caption{Beispiele aus dem Webtest: korrekt erkannte Monstera (links),
fehlklassifizierter Gummibaum (Mitte) und fehlklassifizierte Efeutute (rechts).}
\end{figure}

% ============================================================================
\section{Diskussion und Grenzen}
% ============================================================================
\label{sec:diskussion}

\begin{itemize}[nosep]
  \item \textbf{Validierung $\neq$ echter Test:} Die hohen Genauigkeiten gelten
        auf Daten aus derselben Quelle. Der Webtest (Abschnitt~\ref{sec:webtest})
        zeigt den realistischeren Wert.
  \item \textbf{Laborbilder (PlantVillage):} Die 99\,\% sind durch den sauberen,
        einheitlichen Bildaufbau begünstigt. Reale Feldfotos sind schwieriger.
  \item \textbf{Mögliche Daten-Leckage:} Der „augmentierte“ PlantVillage-Datensatz
        enthält rotierte/gespiegelte Kopien. Liegen Varianten desselben Blattes in
        Training \emph{und} Validierung, wird die Accuracy zu optimistisch
        geschätzt -- ein wichtiger Punkt für die kritische Bewertung.
  \item \textbf{Klassen-Ungleichgewicht:} Bei den Arten reicht die Bildanzahl von
        66 (Yucca) bis 547 (Monstera). Seltene Klassen sind erwartungsgemäß
        schwächer.
  \item \textbf{Kein separater Test-Split:} Für die Zimmerpflanzen wurde mangels
        Daten kein eigener Held-out-Test angelegt; der Webtest ersetzt ihn nur
        teilweise (kleine Stichprobe, nur „gesunde“ Pflanzen).
\end{itemize}

% ============================================================================
\section{Verbesserungsvorschläge}
% ============================================================================
Da es sich um eine Lernarbeit handelt, hier konkrete, gut umsetzbare nächste
Schritte -- vom Verständnis her wertvoller als reines „mehr Accuracy“:

\begin{enumerate}[label=\textbf{\arabic*.}, nosep, itemsep=2pt]
  \item \textbf{Sauberer Test-Split.} Einen festen, nie zum Training/zur
        Modellwahl genutzten Testdatensatz anlegen und nur \emph{einmal} am Ende
        auswerten. Das ist die wichtigste methodische Verbesserung.
  \item \textbf{Daten-Leckage prüfen.} Beim PlantVillage-Datensatz sicherstellen,
        dass augmentierte Varianten desselben Originals nicht gleichzeitig in
        Train und Val landen (Split nach Original-ID statt nach Bild).
  \item \textbf{Grad-CAM einbauen.} Eine Heatmap visualisiert, \emph{worauf} das
        Netz achtet. Das macht Fehlklassifikationen (z.\,B.\ Gummibaum
        $\rightarrow$ Schefflera) erklärbar und ist didaktisch sehr lehrreich.
  \item \textbf{Konfidenz-Schwelle / „unsicher“.} Bei niedriger Konfidenz
        (vgl.\ die 37\,\% beim Pothos) lieber „unsicher“ ausgeben statt einer
        falschen Klasse. Erhöht die Vertrauenswürdigkeit spürbar.
  \item \textbf{Mehr Augmentierung \& Epochen für kleine Datensätze.} Besonders
        das Zustand-Modell (904~Bilder) würde von stärkerer Augmentierung und
        ggf.\ \emph{Early Stopping} profitieren.
  \item \textbf{Reproduzierbarkeit.} Globalen Zufalls-Seed setzen und
        Paketversionen in \texttt{requirements.txt} fixieren (\texttt{pip freeze}),
        damit Läufe exakt wiederholbar sind.
  \item \textbf{Kleiner Code-Hinweis.} \texttt{torch.load} sollte künftig mit
        \texttt{weights\_only=True} aufgerufen werden (PyTorch-Warnung); für
        selbst erzeugte Checkpoints unkritisch, aber gute Praxis.
  \item \textbf{Eigene Feldfotos sammeln.} Ein kleiner, selbst fotografierter
        Testsatz (verschiedene Hintergründe/Licht) zeigt am ehrlichsten, wie gut
        das Modell in der Praxis generalisiert.
\end{enumerate}

% ============================================================================
\section{Fazit}
% ============================================================================
Die Projektarbeit setzt eine vollständige Bildklassifikations-Pipeline mit
Transfer Learning um und deckt alle wesentlichen Schritte ab: Datenaufbereitung,
Zwei-Phasen-Training, Evaluation, Modellvergleich und Inferenz auf neuen Bildern.
Auf den Validierungsdaten werden sehr hohe Werte erreicht (Krankheiten bis
99{,}7\,\%, Arten 91{,}3\,\%); der unabhängige Webtest (8/10) ordnet diese
ehrlich ein und macht den Unterschied zwischen Validierungs- und realer
Leistung sichtbar. Damit ist das eigentliche Lernziel erreicht: ein praktisches
\emph{und} kritisches Verständnis dafür, wie moderne Bildklassifikation
funktioniert und wo ihre Grenzen liegen.

\end{document}
